We now present a series of examples that illustrate the central features of
\Gweek{} and the programming patterns they enable. These examples are chosen not
merely to showcase functionality, but to demonstrate the interaction between the
three FLP effects (choice, logical variables, constraints), the suspension
mechanism, and the narrowing strategy. Each example highlights a different
aspect of the design.

\subsection{Running Functions Backwards}

The paradigmatic example of FLP is `running a function backwards': given a
standard functional definition, using logical variables and constraints to
compute \emph{inverses}. Given the standard list concatenation:
%
\begin{haskellcode}
cat :: [Nat] -> [Nat] -> [Nat]
cat xs ys = case xs of
      [] -> ys
    | (x:xs) -> x : (cat xs ys).
\end{haskellcode}
%
\noindent we can compute the last element of a list by searching for a
decomposition:
%
\begin{haskellcode}
last :: [Nat] -> Nat
last xs = exists ys :: [Nat]. exists y :: Nat.
    cat ys [y] =:= xs. y.
\end{haskellcode}
%
\noindent Evaluating \mintinline{haskell}|last [1,2,3]| introduces logical
variables \mintinline{haskell}|ys| and \mintinline{haskell}|y|, then attempts
to unify \mintinline{haskell}|cat ys [y]| with \mintinline{haskell}|[1,2,3]|.
The unification triggers the evaluation of \mintinline{haskell}|cat ys [y]|,
which pattern-matches on its first argument \mintinline{haskell}|ys|. Since
\mintinline{haskell}|ys| is a logical variable, the machine \emph{narrows}:
it splits into a branch where $\mathit{ys} = \texttt{[]}$ and a branch where
$\mathit{ys} = h : t$ for fresh variables $h$ and $t$. The narrowing continues
recursively until the concatenation matches \mintinline{haskell}|[1,2,3]|,
binding \mintinline{haskell}|y| to \mintinline{haskell}|3|.

This example illustrates the completeness of weak head narrowing: the logical
variable \mintinline{haskell}|ys| is refined \emph{only} because it appears as
the scrutinee of a case expression inside \mintinline{haskell}|cat|. If it were
never scrutinised, it would remain unbound.

\subsection{Permutation Sort}

Sorting can be expressed as a search problem: find a permutation of the input
that is sorted. This demonstrates how nondeterministic choice (\texttt{<>})
combines with the suspension mechanism to enable early pruning.
%
\begin{haskellcode}
insert :: Nat -> [Nat] -> [Nat]
insert x xs = case xs of
    [] -> [x]
  | (z:zs) -> ((x : z : zs) <> (z : insert x zs)).

perm :: [Nat] -> [Nat]
perm xs = case xs of
    [] -> []
  | (z:zs) -> insert z (perm zs).

guard_sorted :: [Nat] -> Nat
guard_sorted xs = case xs of
    [] -> 0
  | (y:ys) -> (case ys of
      [] -> 0
    | (z:zs) -> (let c = leq y z in guard_sorted (z:zs))).

let ys = perm [5, 3, 1, 4, 2] in
let c = guard_sorted ys in
ys.
\end{haskellcode}
%
\noindent The function \mintinline{haskell}|perm| generates all permutations
using choice. The function \mintinline{haskell}|guard_sorted| checks
sortedness, failing on any unsorted pair. The use of \texttt{let} bindings for
both the permutation and the constraint means that both are \emph{suspended}.
The suspension mechanism ensures that the sortedness check is interleaved with
the generation: as each element of the permuted list is demanded, the check
runs on the elements seen so far, pruning unsorted prefixes before the complete
permutation is constructed.

\subsection{N-Queens}

The $n$-queens problem asks for all placements of $n$ queens on an $n \times n$
chessboard such that no two attack each other. This example demonstrates the
importance of \emph{forcing} constraints, as discussed in
\cref{section:suspensions}.
%
\begin{haskellcode}
safe :: Nat -> Nat -> [Nat] -> Nat
safe q d rest = case rest of
    [] -> 0
  | (r:rs) -> case neq (add q d) r of
        Z -> (case neq (add r d) q of
            Z -> safe q (S d) rs
          | S n -> fail)
      | S n -> fail.

queens :: [Nat] -> [Nat] -> [Nat] -> [Nat]
queens skipped rest placed = case rest of
    [] -> (case skipped of
        [] -> placed
      | (y:ys) -> fail)
  | (q:qs) ->
      (case safe q 1 placed of
         Z -> (let remaining = revappend skipped qs in
               case remaining of
                 [] -> (q : placed)
               | (y:ys) -> queens [] remaining (q : placed))
       | S n -> fail)
      <>
      (queens (q : skipped) qs placed).

queens [] [1,2,3,4,5,6,7,8] [].
\end{haskellcode}
%
\noindent The safety check uses \texttt{case} expressions---not \texttt{let}
bindings---to scrutinise the results of \mintinline{haskell}|neq|. This is
essential: a \texttt{let} binding would suspend the check, allowing the machine
to place queens in conflicting positions before discovering the conflict. The
\texttt{case} expression forces the check immediately, pruning invalid branches
at each step. The 8-queens problem finds all 92 solutions in under 200ms.

\subsection{Magic Squares}

A $3 \times 3$ magic square arranges the numbers 1--9 so that each row, column,
and diagonal sums to 15. This example illustrates the dramatic effect of
\emph{interleaving generation and constraint checking}, the second programming
pattern of \cref{section:suspensions}.

A naive approach generates all permutations and checks each:
%
\begin{haskellcode}
let sq = perm [1,2,3,4,5,6,7,8,9] in
exists a :: Nat. exists b :: Nat. (* ... *)
sq =:= [a, b, c, d, e, f, g, h, i].
add3 a b c =:= 15. (* ... *)
sq.
\end{haskellcode}
%
\noindent This must explore all $9! = 362{,}880$ permutations before any
constraint prunes the search, and takes over 30 seconds.

The efficient version uses an incremental selection function that
nondeterministically picks one element from a list:
%
\begin{haskellcode}
pick :: [Nat] -> [Nat] -> [Nat]
pick skipped rest = case rest of
    [] -> fail
  | (x:xs) -> (x : revappend skipped xs)
              <> (pick (x : skipped) xs).
\end{haskellcode}
%
\noindent Elements are picked one at a time, and row-sum constraints are checked
as soon as each row of three is complete. After picking $a$, $b$, $c$, the
constraint \mintinline{haskell}|add3 a b c =:= 15| runs immediately, pruning
the vast majority of branches before $d$, $e$, $f$ are ever chosen. This
version completes in under 250ms---a factor of $120\times$ improvement, arising
purely from the order in which nondeterministic choices and constraints are
combined.

\subsection{Searching over Structures}

Logical variables can be used to search over structured data, demonstrating how
narrowing lazily enumerates the inhabitants of a type.
%
\begin{haskellcode}
exists xs :: [Nat]. length xs =:= 7. sum xs =:= 5. xs.
\end{haskellcode}
%
\noindent The machine introduces a logic variable \mintinline{haskell}|xs| of
type \mintinline{haskell}|[Nat]|. The constraint
\mintinline{haskell}|length xs =:= 7| forces the machine to narrow
\mintinline{haskell}|xs|: the function \mintinline{haskell}|length|
case-analyses its argument, triggering a split into $\texttt{[]}$ and
$h : t$. Only the cons branch can produce a list of length 7, so
\mintinline{haskell}|xs| is progressively instantiated as a 7-element list of
logic variables. The constraint \mintinline{haskell}|sum xs =:= 5| then narrows
each element. The program finds all 462 solutions in under 100ms.

\subsection{Pythagorean Triples}

Logical variables over $\NatT$ can solve simple number-theoretic problems:
%
\begin{haskellcode}
exists a :: Nat. exists b :: Nat. exists c :: Nat.
add (mult a a) (mult b b) =:= (mult c c).
c =:= 5.
(a, (b, c)).
\end{haskellcode}
%
\noindent This searches for Pythagorean triples $(a, b, c)$ with $c = 5$ by
narrowing $a$ and $b$ over the natural numbers and checking the constraint at
each instantiation. The constraint $c \mathbin{\texttt{=:=}} 5$ bounds the
search, as the multiplication $c \cdot c = 25$ limits the range of $a$ and $b$.

\subsection{Constraint Satisfaction}

Constraint satisfaction problems are a natural fit for FLP. The following
program colours the map of Australia with three colours such that no two
adjacent regions share a colour:
%
\begin{haskellcode}
color :: Nat -> Nat
color n = 0 <> 1 <> 2.

let wa = color 0 in let nt = color 0 in
let sa = color 0 in let q  = color 0 in
let nsw = color 0 in let v = color 0 in
let t = color 0 in
let c1 = neq wa nt in let c2 = neq wa sa in
(* ... further adjacency constraints ... *)
[wa, nt, sa, q, nsw, v, t].
\end{haskellcode}
%
\noindent The function \mintinline{haskell}|color| returns one of three colours
nondeterministically. The constraints \mintinline{haskell}|neq| between adjacent
regions are suspended until the result list is constructed, at which point all
variables must be resolved. The program finds all valid colourings.

This example also illustrates the subtlety of the suspension mechanism: the
\mintinline{haskell}|neq| constraints are bound by \texttt{let}, not forced by
\texttt{case}. They are therefore checked only during the drain phase, after
all colours have been chosen. For a map with many constraints, forcing the
constraints earlier---by wrapping them in \texttt{case} expressions---would
provide better pruning.
